\section{Introduction}

Free entropy connects the joint moments of noncommuting variables with
the geometry of their matrix approximations. For a bounded self-adjoint
tuple $X=(X_1,\ldots,X_n)$ in a tracial von Neumann algebra, Voiculescu's
microstates free entropy $\chi(X)$ measures the asymptotic volumes of
matrix tuples whose normalized trace moments approximate those of $X$
\cite{Vchi}. His nonmicrostates free entropy $\chi^*(X)$ instead uses
conjugate variables, the free analogues of classical scores, and integrates
free Fisher information along addition of freely independent semicircular
noise \cite{Vstar}. We give both definitions, with their normalizations,
in Section~2.

The two quantities agree for a single self-adjoint variable, and their
agreement for general tuples was one of the questions in Voiculescu's
unification problem \cite[Sections~3.5 and~3.8]{Vsurvey}.
The finite-entropy formulation asks whether
\[
 \chi(X)>-\infty\quad\Longrightarrow\quad\chi(X)=\chi^*(X).
\]
Guionnet stated this formulation explicitly
\cite[Conjecture~7.1]{Gsurvey}, using the equivalent Gaussian-relative
normalization. The finiteness hypothesis is substantial: as discussed
by Jekel and Pi \cite[Introduction]{JP}, there are nonapproximable laws
with $\chi=-\infty<\chi^*$. Here we construct a strict
gap with both entropies finite.

Biane, Capitaine, and Guionnet proved the general comparison
$\chi\leq\chi^*$ through large-deviation bounds for matrix Brownian
motion \cite{BCG}. Jekel and Pi obtained an elementary proof using
classical entropy and Fisher information, and extended the comparison
to conditional free entropy \cite{JP}. Equality is known in significant
regular convex settings. Dabrowski proved it for a class of strictly
convex potentials with Lipschitz derivatives and noncommutative polynomial
approximations \cite[Theorem~A]{Dab}; Jekel established
entropy and Fisher-information convergence for convex matrix potentials
with appropriate regularity, concentration, and trace-polynomial
approximation \cite[Theorem~7.1]{Jconv}. These results also identify
conditions under which the upper and lower matrix-size limits agree.
Our construction shows that finite entropy alone does not imply equality.

\begin{theorem}\label{thm:main}
There exist an integer $n\geq2$ and a bounded self-adjoint tuple
$X=(X_1,\ldots,X_n)$ in a von Neumann algebra with faithful normal
tracial state such that
\[
 -\infty<\chi(X)\leq\chi^*(X)-\tfrac12<\infty.
\]
Here $\chi$ is the original microstates entropy using an operator-norm
cutoff and a limsup over matrix sizes.
\end{theorem}

Thus Theorem~\ref{thm:main} answers the finite-microstates-entropy
equality question negatively. The construction uses a large but fixed number of
variables. It does not address equality for every smaller prescribed
number of variables.

The proof compares two noise scales. Start with
$C=(B_1,\ldots,B_p,Y)$, where $B_1,\ldots,B_p$ are free variance-one
semicircular variables and $Y$, in a separate tensor factor, is uniformly
distributed on $p$ equally spaced points in $[-1,1]$. Thus $Y$ commutes
with every $B_i$. Let $S$ be a standard semicircular $(p+1)$-tuple free
from $C$. A small perturbation $X=C+\sqrt h S$ has finite microstates
entropy, while its last coordinate still determines spectral blocks
containing much of the energy of the first $p$ coordinates.

After an additional noise time $t$, we obtain two different estimates.
The analytic estimate uses Speicher's noncrossing cumulant calculus
\cite{Speicher,NS} and orthogonality of semicircular Wick polynomials.
Proposition~\ref{prop:fisher} bounds the Fisher deficit by squared
$\ell^2$ norms of the higher cumulant arrays. The sparsity of the signal's
cumulants then makes the nonmicrostates entropy of
$C+\sqrt{h+t}\,S$ close to the entropy of the semicircular family with
the same coordinate variances.
The matrix estimate shows that uniform laws on a selected sequence of
increasingly accurate cutoff microstate sets of $X$, evolved by independent
GUE noise of variance $t$, retain detectable simultaneous block structure.
This event has probability bounded away from zero uniformly along
the selected matrix sizes for those ensembles, and probability at most
$e^{-c d^2}$ for size-$d$ independent GUE matrices with the coordinate
variances of the limiting law, for a fixed $c>0$.
Relative-entropy data processing therefore gives a deficit bounded away
from zero in normalized classical entropy. The use of commutation constraints and projection
geometry to control microstate volumes has antecedents in Ge's argument
for prime factors \cite[Lemma~2 and its proof outline]{Ge}.
The unitary-net estimate below uses the elementary volumetric covering
method discussed by Szarek \cite{Szarek}.
The event estimated here concerns the persistence of block energy after
positive Gaussian time. The analytic entropy deficit becomes small when
$t/\sqrt p\to\infty$, while the Gaussian block estimate still beats the
unitary covering cost when $t\log p/p\to0$. The choice $t=p^{3/4}$ satisfies
both conditions as $p\to\infty$. We first choose one sufficiently large
finite $p$ and then $h$; all matrix limits are taken with these parameters fixed.

The transfer between the scales uses the classical-to-free Fisher
comparison of Jekel and Pi \cite[Proposition~6.1]{JP}.
Lemma~\ref{lem:flow} proves the needed finite-time statement while retaining
the entropies of this same sequence of terminal laws. The limsup of those
normalized entropies need not equal the microstates entropy of the limiting
law. Keeping this sequence in the comparison permits the positive-time
structural deficit to produce the strict gap for $X$.

The companion paper \emph{An isomorphism of the free group factors}
\cite[Corollary~7.1]{FreeGroupFactors} proves that each of the four free
entropy dimensions $\delta,\delta_0,\delta^*,\delta^\star$ attains every
integer value $n\geq2$ on suitable finite self-adjoint $W^*$-generating
tuples of $L(\mathbb F_2)$. That statement varies the generating tuple of
a fixed algebra; Theorem~\ref{thm:main} compares $\chi$ and $\chi^*$ on one
tuple, with both values finite.

Section~2 fixes the conventions and establishes the matrix convergence
facts, starting from Gaussian asymptotic freeness \cite{Vgue,NS}.
Section~3 proves the cumulant/Fisher estimate. Section~4 constructs the
signal, fixes the two noise scales, and proves $\chi(X)>-\infty$.
Section~5 proves the specified-ensemble entropy-production
lemma. Section~6 establishes the block event and combines the estimates
to prove Theorem~\ref{thm:main}.
